\begin{center}\textbf{Abstract}\end{center}
\small
An AI coding agent can hand some of its work to a cheaper model. Whether that saves money in a long
conversation is hard to estimate, because providers cache the conversation separately for each model and
a switch throws that cache away. We stopped estimating and measured. For each of \NScenarios\ scripted
\MinTurns--\MaxTurns-turn sessions, the plain host model, the shipped router (which re-chooses its model at
the start of each turn), a ``sticky'' router that chooses a model once and never switches, and plain Sonnet
(plus, on a \AAPct\,\% subsample, a second plain copy) were launched within seconds of each other from one
frozen workspace, each with its own cache key, twice on each of two host models: \NSessions\ sessions in
all. Hypotheses, thresholds, split and decision rule were preregistered; the estimator was written after the
data were collected. On the \NTest\ held-out scenarios, with Fable~5.1 as host at our price table, sticky
routing cost \CfFableStickyPairTwo\X\ as much as the plain host (95\,\% confidence interval
\CfFableStickyPairCITwo), with turn-pass fraction inside the preregistered margin; plain Sonnet cost about
the same (\CfFableSonnetPairTwo\X). With Opus~5.5 as host, every routing arm cost more
(\CfOpusStickyPairTwo\X--\CfOpusSonnetPairTwo\X; \CfOpusStickyArmTwo\X--\CfOpusSonnetArmTwo\X\ without the
quality filter); that result depends on Opus's cheap cache reads and would reverse if they cost about
\BEStickyPair\ instead of \PriceOpusRead\ dollars per million tokens. The sticky policy cost less than the
shipped one on both hosts. A savings model fitted on the training scenarios passed its preregistered pooled
check on the held-out ones. The measurement cost \$\SpendLedger.
\normalsize
